\section{Conclusion}
\label{sec_conclusion_performance}

The results presented in this chapter demonstrate that the integration of GPU acceleration and Docker-based deployment in DataLex significantly enhances its performance compared to traditional SIEM solutions. The improvements observed in inference time, retrieval latency, memory utilization, throughput, and token generation rate clearly indicate that leveraging high-performance hardware and optimized architectures allows for real-time security monitoring at scale. The multi-threaded capabilities of Suricata further strengthen intrusion detection efficiency, surpassing Snort in terms of packet processing speed and log integration. Additionally, by adopting a containerized SIEM model, DataLex achieves greater scalability, resource efficiency, and cost-effectiveness compared to traditional solutions such as Splunk and Fortinet, which rely on monolithic or appliance-based deployments. These advancements not only ensure faster security analytics and incident response but also provide a future-proof foundation for AI-driven threat intelligence in enterprise environments.

As security threats continue to evolve, the findings of this study confirm that AI-accelerated, Docker-based SIEM solutions like DataLex are well-positioned to meet the growing demands of real-time cybersecurity monitoring and compliance enforcement.
